\section{A Zoo of Syntaxes}\label{implicit-presentations}

\subsection{Implicit presentations}

In this section, we introduce the more implicit presentations of type theory, which we will prove equivalent to \ThyGAT.

First, we describe \ThyTarski, which has implicit substitutions and exactly the same rules as the GAT 
of section $2$. 
A main goal of this paper is to relate \ThyTarski to the implicit formulation \textit{à la Russell} of cumulative universes denoted with \ThyRussell.
The latter conflates types and terms, with implicit inclusion between universes. 
This system cannot be formulated algebraically, as this would involve adding a sort equation and a stratification of the sort of types $\Type_n$ which we would need to be included in each other.

Both \ThyTarski and \ThyRussell have dependent products, lambda abstraction and application, described in the same way, with the
following rules.


\begin{figure}[H]
    \centering
    \setlength{\fboxsep}{10pt}
    \fbox{% 
        \begin{minipage}{0.95\textwidth}
\begin{mathparsmall}
  \ebrule{
    \hypo{\Gamma \vdash A \;\type}
    \infer1{\Gamma.A \vdash v_0 : A\p}
  }
  \and
  \ebrule{
    \hypo{\Gamma \vdash v_n:A}
    \hypo{\Gamma \vdash B \;\type}
    \infer2{\Gamma.B \vdash v_{n+1} : A\p}
  }
  \and
  \ebrule{
    \hypo{\Gamma \vdash A \;\type}
    \hypo{\Gamma.A \vdash B \;\type}
    \infer2{\Gamma \vdash \Pi(A, B) \;\type}
  }
  \and
  \ebrule{
    \hypo{\Gamma.A \vdash b: B}
    \hypo{\Gamma.A \vdash B \;\type}
    \hypo{\Gamma \vdash A \;\type}
    \infer3{\Gamma \vdash \lambda(A,B,b) : \Pi (A, B)}
  }
  \and
  \ebrule{
    \hypo{\Gamma \vdash c:\Pi(A, B)}
    \hypo{\Gamma \vdash a: A}
    \hypo{\Gamma.A \vdash B \;\type}
    \hypo{\Gamma \vdash A \;\type}
    \infer4{\Gamma \vdash \app(A,B,c,a) : B[a]}
  }
  \and
  \ebrule{
    \hypo{\Gamma.A \vdash b: B}
    \hypo{\Gamma \vdash a:A}
    \infer2[($\beta$-rule)]{\Gamma \vdash \app(A,B,\lambda(A,B,b),a) = b[a] : B[a]}
  }
  \and
  \ebrule{
    \hypo{\Gamma \vdash f: \Pi \; A \; B}
    \infer1[($\eta$-rule)]{\Gamma \vdash f = \lambda(A,B,\app(A\p,B\p^+,f\p, v_0)) : \Pi(A, B)}
  }
\end{mathparsmall}
\end{minipage}%
    }
    \caption{Logical structure for \ThyTarski and \ThyRussell}
    \label{fig:rulestyping}
\end{figure}

Then, the operations of \ThyGAT can be translated as formation and introduction rules for Tarski-style universes and the theory
\ThyTarski. Namely, we need codes for products and universes, lifts and decoding. We describe the rules in the following:

\begin{figure}[H]
    \centering
    \setlength{\fboxsep}{10pt}
    \fbox{% 
        \begin{minipage}{0.95\textwidth}
\begin{mathparsmall}
  \ebrule{
    \hypo{}
    \infer1{\Gamma \vdash \U_n \;\type}
  }
  \and
  \ebrule{
    \hypo{\Gamma \vdash a: \U_n}
    \hypo{\Gamma.T_n(a) \vdash b: \U_n}
    \hypo{\Gamma \vdash T_n(a) \;\type}
    \infer3{\Gamma \vdash \Pi^n \; a \; b: \U_n}
  }
  \and
  \ebrule{
    \hypo{\Gamma \vdash a:\U_n}
    \infer1{\Gamma \vdash T_n(a) \;\type}
  }
  \and
  \ebrule{
    \hypo{}
    \infer1[($m<n$)]{\Gamma \vdash \U_m^n: \U_n}
  }
  \and
  \ebrule{
    \hypo{\Gamma \vdash a: \U_m}
    \infer1[($m\leq n$)]{\Gamma \vdash t_m^n(a): \U_n}
  }
\end{mathparsmall}
\end{minipage}%
    }
    \caption{Universe rules of \ThyTarski}
    \label{fig:rulesconv}
\end{figure}

The equations of \ThyGAT can be translated as conversion rules for \ThyTarski (Figure \ref{fig:rulesconv}). Namely,
we require the decoding and lifting operation to be strict, and natural.

\begin{figure}[H]
    \centering
    \setlength{\fboxsep}{10pt}
    \fbox{% 
        \begin{minipage}{0.95\textwidth}
\begin{mathparsmall}
  \ebrule{
    \hypo{\Gamma \vdash a: \U_n}
    \hypo{\Gamma.T_n(a) \vdash b: \U_n}
    \infer2{\Gamma \vdash T_n(\Pi^n \; a \; b) = \Pi(T_n(a),T_n(b))}
  }
  \and
  \ebrule{
    \hypo{\Gamma \vdash a: \U_n}
    \hypo{\Gamma.T_n(a) \vdash b: \U_n}
    \infer2[($n\leq p$)]{\Gamma \vdash t_n^p(\Pi^n \; a \; b) = \Pi^p(t_n^p(a),t_n^p(b)) : \U_p}
  }
  \and
  \ebrule{
    \hypo{}
    \infer1[$m<n$]{\Gamma \vdash T_n(\U_m^n) = \U_m}
  }
  \and
  \ebrule{
    \hypo{}
    \infer1[($n\leq p$)]{\Gamma \vdash t_n^p(\U_m^n) = \U_m^p : \U_p}
  }
  \and
  \ebrule{
    \hypo{\Gamma \vdash a : \U_m}
    \infer1[($m\leq n$)]{\Gamma \vdash T_n(t_m^n(a)) = T_m(a)}
  }
  \and
  \ebrule{
    \hypo{\Gamma \vdash a : \U_m}
    \infer1[($m\leq n\leq p$)]{\Gamma \vdash t_n^p(t_m^n(a)) = t_m^p(a): \U_p}
  }
\end{mathparsmall}
\end{minipage}%
    }
    \caption{Conversion rules of \ThyTarski}
    \label{fig:rulesconv}
\end{figure}

Now, for the theory \ThyRussell (Russell universes), types and terms are conflated, and thus
we should be able to deduce $\Gamma \vdash A \; \type$ from $\Gamma \vdash A : \U_n$. Conversely,
the intuition is that whenever $\Gamma \vdash A \; \type$ then we should have  : \[\exists n \in \bN, \ \Gamma \vdash A : \U_n\]

Instead of having codes for each of the type formers, we now close $\U_n$ under the type connectives. Note that, as we
consider \textit{cumulative} universes, we do not take the suprema of levels for the dependent product. Cumulativity itself
is expressed by asking for the proper inclusion of universes : if $\Gamma \vdash A : \U_n$ then $\Gamma \vdash A : \U_p$ for $p > n$. 

\begin{figure}[H]
    \centering
    \setlength{\fboxsep}{10pt}
    \fbox{% 
        \begin{minipage}{0.95\textwidth}
\begin{mathparsmall}
\ebrule{
    \hypo{}
    \infer1{\Gamma \vdash \U_n \;\type}
  }
  \and
  \ebrule{
    \hypo{\Gamma \vdash A:\U_n}
    \infer1{\Gamma \vdash A \;\type}
  }
  \and
  \ebrule{
    \hypo{\Gamma \vdash A: \U_n}
    \hypo{\Gamma.A \vdash B: \U_n}
    \infer2{\Gamma \vdash \Pi(A, B): \U_n}
  }
  \and
  \ebrule{
    \hypo{}
    \infer1[($m<n$)]{\Gamma \vdash \U_m: \U_n}
  }
  \and
  \ebrule{
    \hypo{\Gamma \vdash A: \U_n}
    \infer1[($n<p$)]{\Gamma \vdash A: \U_p}
  }
\end{mathparsmall}
\end{minipage}%
    }
    \caption{Universe rules of \ThyRussell}
    \label{fig:rulesconv}
\end{figure}

Finally, we also consider a system \ThyPTS, which removes redundant type annotations in functional application
and lambda abstraction from \ThyRussell. Thus, $\app(A,B,c,a)$ and $\lambda(A,B,b)$ become $c \; a$ and $\lambda_{x:A} b$.
This system is closer to the practical syntax of type theory, inherited from the lambda calculus,
and often referred to as the "pure type system" presentation (PTS).


For a full and extensive description of the systems, please refer to appendix \ref{appendix:rules}.

\subsection{Forgetful maps}

Now, we relate the previously defined theories, by exhibiting forgetful translations from the more explicit presentations to the less explicit ones.

First, we define a morphism of theories between \ThyGAT and \ThyTarski which makes substitutions implicit.
To do this, we map substitutions in \ThyGAT to a list of elementary substitutions in \ThyTarski, as there is no composition in this system. Then, when translating the action of a substitution
on a term or a type, we can unfold that list to get a type obtained by applying all those substitutions successively.

\begin{definition}
    We define $\fold(.)$ inductively by: 

    \noindent
    \begin{minipage}{0.40\linewidth}
    \begin{align*}
        \fold(\q) &= v_0\\
        \fold([a]) &= [\fold(a)]\\
        \fold(A \sigma) &= (...(\fold(A)\sigma_1)...)...\sigma_n\\
        \fold(a \sigma) &= (...(\fold(a)\sigma_1)...)...\sigma_n\\ 
    \end{align*}
    \end{minipage}
    \hfill
    \begin{minipage}{0.40\linewidth}
    \begin{align*}
        \fold(\p) &= \p\\
        \fold(\Id) &= []\\ 
        \fold(\sigma\delta) &= \fold(\sigma) @ \fold(\delta)\\
        \fold(\sigma^+) &= (\sigma_1^+,..., \sigma_n^+)
    \end{align*}
    \end{minipage}
\end{definition}

Then, we define the following erasure operation, which erases the lift and decoding operation of 
Tarski-style universes, and sends codes to the corresponding type formers. Concretely, this operation
forgets about all the explicit coercions given by the presentation à la Tarski, and leaves everything implicit.

\begin{definition}[\formline{Erasure.v}{10}{Erasurefunction} Erasure]
    We define the erasure operation $|.|$ inductively by:

    \noindent
    \begin{minipage}{0.40\linewidth}
    \begin{align*}
        |\U_n| &= \U_n\\
        |\U_m^n| &= \U_m\\
        |\app(A,B,c,a)| &= \app(|A|,|B|,|c|,|a|)\\
        |\lambda(A,B,b)| &= \lambda(|A|,|B|,|b|)\\
        |v_i| &= v_i\\
    \end{align*}
    \end{minipage}
    \hfill
    \begin{minipage}{0.40\linewidth}
    \begin{align*}
        |T_n(a)| &= |a|\\
        |t_m^n(a)| &= |a|\\
        |\Pi^n \; a \; b| &= \Pi(|a|, |b|)\\
        |\Pi(A, B)| &= \Pi(|A|, |B|)
    \end{align*}
    \end{minipage}
\end{definition}

Finally, we define the map $\strip$  which erases type annotations from a term of the system \ThyRussell, and gives a term of \ThyPTS.
This is analogous to the one defined by \citet{Streicher}.

\begin{definition}
    We define $\strip$ inductively by :
    \begin{align*}
        \strip(\lambda (A,B,b)) &= \lambda(\strip(A), \strip(b))\\ 
        \strip(\app (A,B,c,a)) &= \strip(c) \; \strip(a)\\
        \strip(v_i) &= v_i 
        \strip(\Pi(A,B)) &= \Pi(\strip(A), \strip(B))\\
        \strip(\U_n) &= \U_n
    \end{align*}
\end{definition}

In the following, we will prove lemmas of the following shape, following \citet{Streicher}:\\

\textbf{Lemma.} If $\Gamma \vdash t,t' : A$ in $\mathcal{T}$ with $|t| = |t'|$ in $\mathcal{T}'$, then we have : $\Gamma \vdash t = t' : A$ in $\mathcal{T}$.\\

\begin{remark}As mentioned in the introduction, this relates to elaboration in proof assistants. 
Concretely, this lemma means that if two terms are syntactically the same in the implicit user syntax, then they should have the same
computational behaviour in the core system, and thus should be convertible. This is to be compared with a simulation
result, and is often referred to as a \emph{coherence} lemma.
\end{remark}

In the next sections, we will prove such lemmas for each of our theories. With those, we will be able to lift a derivation from an implicit theory to an explicit one.
The situation is summarized in Figure \ref{hierarchy-zoo}, where the strong arrows are the known forgetful maps, and the dotted arrows are the sections we will be constructing.

\textcolor{red}{TODO: precise what kind of equivalence this is. What does this means in terms of models, syntax etc...}

\begin{figure}[!ht]
\centering
\resizebox{\textwidth}{!}{%
\begin{circuitikz}
\tikzstyle{every node}=[font=\LARGE]

\draw  (-2.5,18.75) rectangle  node {\LARGE GAT (\ThyGAT)} (2.5,16.25);
\draw  (5,18.75) rectangle  node {\LARGE Tarski (\ThyTarski)} (10,16.25);
\draw  (12.5,18.75) rectangle  node {\LARGE Russell (\ThyRussell)} (17.5,16.25);
\draw  (20,18.75) rectangle  node {\LARGE PTS (\ThyPTS)} (25,16.25);






\draw [->, >=Stealth] (2.5,18.25) -- (5,18.25)node[pos=0.5, fill=white]{$\fold$};
\draw [->, >=Stealth] (10,18.25) -- (12.5,18.25)node[pos=0.5, fill=white]{$|.|$};
\draw [->, >=Stealth] (17.5,18.25) -- (20,18.25)node[pos=0.5, fill=white]{$\strip$};
\draw [->, >=Stealth, dashed] (5,16.75) -- (2.5,16.75);
\draw [->, >=Stealth, dashed] (12.5,16.75) -- (10,16.75);
\draw [->, >=Stealth, dashed] (20,16.75) -- (17.5,16.75);
\end{circuitikz}
}%

\caption{The hierarchy of theories}
\label{hierarchy-zoo}
\end{figure}

As mentioned in the introduction, to prove our equivalence result, we rely on the definitional inversion properties for \textit{only} the explicit system
\ThyGAT proved in section \ref{GAT}, and we propagate those along our equivalence results.
In the following sections \ref{sectionT0T1}, \ref{sectionT1T2}, and \ref{sectionT2T2'} we show that all the systems \ThyGAT, \ThyTarski, \ThyRussell and \ThyPTS are equivalent.
